\documentclass[12pt]{article}
\usepackage[brazil]{babel}
\usepackage[utf8]{inputenc}
\usepackage[T1]{fontenc}
\usepackage{amsmath, amssymb}
\usepackage{geometry}
\usepackage{graphicx}
\usepackage{tikz}
\usetikzlibrary{shapes.gates.logic.US, positioning}
\usepackage{xcolor}
\usepackage{enumitem}

\definecolor{maincolor}{RGB}{0,100,50}
\usepackage{fancyhdr}
\definecolor{forestgreen}{RGB}{0,100,50}
\usepackage[most]{tcolorbox}

\newtcolorbox{block}[1]{%
  colback=green!5,
  colframe=maincolor,
  coltitle=white, 
  title=#1,
  fonttitle=\bfseries,
  boxrule=1.5pt,
  arc=3mm,
  left=6pt,
  right=6pt,
  top=6pt,
  bottom=6pt
}

\newtcolorbox{codeblock}{%
  colback=black!2,
  colframe=gray!30,
  arc=1mm,
  boxrule=0.5pt,
  left=5pt,
  right=5pt,
  top=5pt,
  bottom=5pt,
  fontupper=\small
}

\geometry{margin=2.5cm, headheight=53pt, top=2.5cm}
\pagestyle{fancy}
\fancyhf{}

\renewcommand{\headrulewidth}{2pt}
\renewcommand{\footrulewidth}{2pt}

\renewcommand{\headrule}{\hbox to\headwidth{%
  \color{forestgreen}\leaders\hrule height \headrulewidth\hfill}}

\renewcommand{\footrule}{\hbox to\headwidth{%
  \color{forestgreen}\leaders\hrule height \footrulewidth\hfill}}

\fancyhead[L]{%
  \includegraphics[height=1.2cm]{Universidade_de_Rio_Verde.jpg}
}

\fancyhead[R]{%
  \raggedleft
  \textcolor{forestgreen}{%
    \footnotesize 
    \textbf{Lista de Exercícios — Álgebra Booleana}
    Me.~Sergio Souza Novak\\
    \textit{Lógica para Engenharia de Software}
  }
}

\fancyfoot[C]{%
  \textcolor{forestgreen}{\thepage}
}

\begin{document}


\begin{block}{Instruções}
  RESPONDA EXCLUSIVAMENTE A CANETA, RESPOSTAS DIGITADAS NÃO SERÃO ACEITAS, TIRE FOTO DAS RESPOSTAS ESCRITAS AO PRÓPRIO PUNHO E CONFECCIONE UM PDF COM AS RESPOSTAS.
Aplique as propriedades da Álgebra Booleana (Identidade, Nulidade, Idempotência, Complemento, Absorção e De Morgan) para simplificar as expressões abaixo. Para as questões de programação, apresente o código final refatorado. Para os circuitos, apresente a expressão simplificada e o diagrama reduzido.
\end{block}


\begin{enumerate}[label=\textbf{\arabic*.}]

% Ex 1
\item Simplifique o seguinte bloco condicional, eliminando redundâncias óbvias:
\begin{codeblock}
\begin{verbatim}
if (usuarioAtivo == true && processamentoConcluido == true) {
    gerarRelatorio();
}
\end{verbatim}
\end{codeblock}

\vspace{0.1cm}


% Ex 2
\item Aplique a \textbf{Lei da Absorção} para simplificar a lógica abaixo:
\begin{codeblock}
\begin{verbatim}
if (isAdmin || (isAdmin && temPermissaoEspecial)) {
    acessoTotal();
}
\end{verbatim}
\end{codeblock}

\vspace{0.1cm}


% Ex 3
\item Utilize o \textbf{Teorema de De Morgan} para simplificar a consulta SQL:

\begin{codeblock}
\begin{verbatim}
SELECT * 
FROM itens
WHERE NOT (itemValido = 1 AND itemExpirado = 0);
\end{verbatim}
\end{codeblock}


\vspace{0.1cm}


% Ex 4
\item Refatore o código abaixo aplicando a \textbf{Idempotência}:
\begin{codeblock}
\begin{verbatim}
if ((logado && ativo) || (logado && ativo)) {
    dashboard();
}
\end{verbatim}
\end{codeblock}

\vspace{0.1cm}


% Ex 5
\item Simplifique a expressão abaixo que utiliza variáveis complementares:
\begin{codeblock}
\begin{verbatim}
if ((podeEditar && usuarioDono) || (podeEditar && !usuarioDono)) {
    salvarAlteracoes();
}
\end{verbatim}
\end{codeblock}

\vspace{0.1cm}



% Ex 6
\item Reduza a complexidade aplicando \textbf{Dupla Negação} e \textbf{De Morgan}:
\begin{codeblock}
\begin{verbatim}
if (!(!conexaoEstavel || !servidorOnline)) {
    iniciarSincronizacao();
}
\end{verbatim}
\end{codeblock}

\vspace{0.1cm}


% Ex 7
\item Identifique a redundância na verificação de limites:
\begin{codeblock}
\begin{verbatim}
if (estoque > 0 && (estoque >= 0 || pedidoUrgente)) {
    separarParaEnvio();
}
\end{verbatim}
\end{codeblock}

\vspace{0.1cm}

% Ex 8
\item Simplifique a função em C abaixo:

\begin{codeblock}
\begin{verbatim}
int deveProcessar(int A, int B) {
    if ((A && B) || (A && !B) || (!A && B)) {
        return 1;
    }
    return 0;
}
\end{verbatim}
\end{codeblock}

\vspace{0.2cm}

\vspace{0.1cm}


% Ex 9
\item Aplique a simplificação para o caso de "Identidade" em negações:
\begin{codeblock}
\begin{verbatim}
if (!(estaChovendo || estaNublado) || (estaChovendo && !estaNublado)) {
    abrirJanelas();
}
\end{verbatim}
\end{codeblock}

\vspace{0.1cm}


% Ex 10
\item Refatore a condição abaixo fatorando os termos comuns:
\begin{codeblock}
\begin{verbatim}
if ((sensorA > 50 && sensorB < 20) || (sensorA > 50 && !(sensorB < 20))) {
    atuaControle();
}
\end{verbatim}
\end{codeblock}

\vspace{0.1cm}

% Ex 11
\item Em uma universidade, uma \textbf{catraca eletrônica} libera a entrada de um aluno com base em duas condições:
\begin{itemize}
\item \texttt{bloqueado}: indica se o aluno está com acesso bloqueado
\item \texttt{conectado}: indica se o sistema conseguiu validar os dados do aluno
\end{itemize}

A lógica atual está definida da seguinte forma:

\begin{codeblock}
\begin{verbatim}
if (!(!bloqueado || !conectado) || !(bloqueado || !conectado)) {
    liberarEntrada();
}
\end{verbatim}
\end{codeblock}

\vspace{0.2cm}

\textbf{Simplifique a expressão lógica acima.}

\vspace{0.1cm}

% Ex 12
\item Reduza a expressão de 3 variáveis para a sua forma mínima:
\begin{codeblock}
\begin{verbatim}
if ((A && B && C) || (A && B && !C) || (A && !B)) {
    executarTarefa();
}
\end{verbatim}
\end{codeblock}

\vspace{0.1cm}

% Ex 13
\item Aplique as propriedades de simplificação para remover a redundância:
\begin{codeblock}
\begin{verbatim}
if ((isAdmin && canEdit) || (!isAdmin && isOwner) || (canEdit && isOwner)) {
    modificarArquivo();
}
\end{verbatim}
\end{codeblock}

\vspace{0.1cm}

% Ex 14
\item Simplifique e identifique o comportamento lógico desta condição:
\begin{codeblock}
\begin{verbatim}
if (!(permiteAcesso || !(permiteAcesso && temToken))) {
    loginError();
}
\end{verbatim}
\end{codeblock}

\vspace{0.1cm}
% Ex 15
\item Em uma aplicação \textbf{Qt (C++)}, um botão executa uma ação ao ser clicado com base nas seguintes condições:
\begin{itemize}
\item \texttt{isValid}: dados válidos
\item \texttt{isCached}: resultado em cache
\item \texttt{hasNetwork}: conexão com a rede disponível
\item \texttt{isAuthorized}: usuário autenticado
\end{itemize}

O código do slot é dado por:

\begin{codeblock}
\begin{verbatim}
void MainWindow::on_pushButton_clicked() {
    if ( ((isValid && hasNetwork) || isCached) &&
         ((!isValid || isCached) && (isAuthorized || isCached)) ) {
        exibirResultado();
    }
}
\end{verbatim}
\end{codeblock}

\vspace{0.2cm}

\textbf{Simplifique a lógica condicional acima.}

\vspace{0.3cm}

% Ex 16
\item Imagine o sistema de controle de resfriamento de uma \textbf{usina nuclear}, onde o alarme só deve ser ignorado (\texttt{false}) se as condições de segurança forem perfeitas. Caso contrário, o alarme dispara (\texttt{true}).

As variáveis são:
\begin{itemize}
\item \texttt{A}: temperatura acima do limite
\item \texttt{B}: pressão interna estável (OK)
\item \texttt{C}: bomba de resfriamento ligada
\item \texttt{D}: sensor de radiação normal (OK)
\end{itemize}

O sistema utiliza a seguinte lógica para decidir se o alarme dispara:
\begin{codeblock}
\begin{verbatim}
bool alarme = !( (A && !B) || (C && !D) ) && ( (A || C) && !(B && D) );
\end{verbatim}
\end{codeblock}

\textbf{Simplifique a expressão lógica acima para encontrar a condição mínima de disparo do alarme.}

\vspace{0.5cm}


% Ex 17
\item Simplifique o circuito de identidade $F = A \cdot A$ e desenhe o equivalente reduzido.

\vspace{0.1cm}
\begin{center}
\begin{tikzpicture}[scale=0.8]
  \node[and gate US, draw, thick] (and) at (0,0) {};
  \draw (and.input 1) -- ++(-1,0) node(pt){};
  \draw (pt.center) node[circle,fill,inner sep=1pt] {} |- (and.input 2);
  \draw (pt.center) -- ++(-0.5,0) node[left] {$A$};
  \draw (and.output) -- ++(0.5,0) node[right] {$F$};
\end{tikzpicture}
\end{center}


% Ex 18
\item Analise o circuito $F = A + \overline{A}$ e determine seu valor constante.
\begin{center}
\begin{tikzpicture}[scale=0.8]
  \node[or gate US, draw, thick] (or) at (2,0) {};
  \node[not gate US, draw, thick] (not) at (0,-0.25) {};
  \draw (or.input 1) -- ++(-2,0) node(pt){};
  \draw (pt.center) -- ++(-0.5,0) node[left] {$A$};
  \draw (pt.center) node[circle,fill,inner sep=1pt] {} |- (not.input);
  \draw (not.output) |- (or.input 2);
  \draw (or.output) -- ++(0.5,0) node[right] {$F$};
\end{tikzpicture}
\end{center}


% Ex 19
\item Aplique a \textbf{Lei da Absorção} no circuito $F = AB + A$.
\begin{center}
\begin{tikzpicture}[scale=0.8]
  \node[and gate US, draw, thick] (and) at (0,0) {};
  \node[or gate US, draw, thick] (or) at (2,0.5) {};
  \draw (and.input 1) -- ++(-1.5,0) node(pt){};
  \draw (pt.center) -- ++(-0.5,0) node[left] {$A$};
  \draw (and.input 2) -- ++(-1,0) node[left] {$B$};
  \draw (pt.center) node[circle,fill,inner sep=1pt] {} |- (or.input 1);
  \draw (and.output) |- (or.input 2);
  \draw (or.output) -- ++(0.5,0) node[right] {$F$};
\end{tikzpicture}
\end{center}


% Ex 20
\item Simplifique a expressão do circuito: $F = (A+B)(A+C)$.
\begin{center}
\begin{tikzpicture}[scale=0.8]
  \node[or gate US, draw, thick] (or1) at (0,0.8) {};
  \node[or gate US, draw, thick] (or2) at (0,-0.8) {};
  \node[and gate US, draw, thick] (and) at (3,0) {};
  \draw (or1.input 1) -- ++(-1.5,0) node(ptA){};
  \node[left] at (ptA) {$A$};
  \draw (or1.input 2) -- ++(-1,0) node[left] {$B$};
  \draw (or2.input 2) -- ++(-1,0) node[left] {$C$};
  \draw (ptA.center) node[circle,fill,inner sep=1pt] {} |- (or2.input 1);
  \draw (or1.output) |- (and.input 1);
  \draw (or2.output) |- (and.input 2);
  \draw (and.output) -- ++(0.5,0) node[right] {$F$};
\end{tikzpicture}
\end{center}



% Ex 21
\item Demonstre a equivalência de De Morgan para a porta NAND: $F = \overline{A \cdot B}$.
\begin{center}
\begin{tikzpicture}[scale=0.8]
  \node[nand gate US, draw, thick] (nand) at (0,0) {};
  \draw (nand.input 1) -- ++(-1,0) node[left] {$A$};
  \draw (nand.input 2) -- ++(-1,0) node[left] {$B$};
  \draw (nand.output) -- ++(1,0) node[right] {$F = \overline{A \cdot B}$};
\end{tikzpicture}
\end{center}


% Ex 22
\item Reduza o circuito $F = AB + \overline{A}B$ fatorando o termo comum.
\begin{center}
\begin{tikzpicture}[scale=0.8]
  \node[and gate US, draw, thick] (and1) at (0,1) {};
  \node[and gate US, draw, thick] (and2) at (0,-0.5) {};
  \node[not gate US, draw, thick] (not) at (-1.5, -0.25) {};
  \node[or gate US, draw, thick] (or) at (2.5, 0.25) {};
  \draw (and1.input 1) -- ++(-2.5,0) node(ptA){}; \node[left] at (ptA) {$A$};
  \draw (and1.input 2) -- ++(-1.5,0) node(ptB){}; \node[left] at (ptB) {$B$};
  \draw (ptA.center) node[circle,fill,inner sep=1pt] {} |- (not.input);
  \draw (ptB.center) node[circle,fill,inner sep=1pt] {} |- (and2.input 2);
  \draw (not.output) |- (and2.input 1);
  \draw (and1.output) |- (or.input 1);
  \draw (and2.output) |- (or.input 2);
  \draw (or.output) -- ++(0.5,0) node[right] {$F$};
\end{tikzpicture}
\end{center}


% Ex 23
\item Simplifique o circuito de absorção complexa: $F = A(\overline{A} + B)$.




% Ex 24
\item Utilize De Morgan para simplificar: $F = \overline{\overline{A} + \overline{B}}$.
\begin{center}
\begin{tikzpicture}[scale=0.8]
  \node[not gate US, draw, thick] (not1) at (0, 1) {};
  \node[not gate US, draw, thick] (not2) at (0, -1) {};
  \node[nor gate US, draw, thick] (nor) at (2, 0) {};
  \draw (not1.input) -- ++(-0.5,0) node[left] {$A$};
  \draw (not2.input) -- ++(-0.5,0) node[left] {$B$};
  \draw (not1.output) |- (nor.input 1);
  \draw (not2.output) |- (nor.input 2);
  \draw (nor.output) -- ++(1,0) node[right] {$F$};
\end{tikzpicture}
\end{center}


% Ex 25
\item Fatore a expressão do circuito de 3 variáveis: $F = ABC + AB\overline{C}$.
\begin{center}
\begin{tikzpicture}[scale=0.7]
  \node[and gate US, draw, thick, logic gate inputs=nnn] (and1) at (0, 2) {};
  \node[and gate US, draw, thick, logic gate inputs=nnn] (and2) at (0, 0) {};
  \node[not gate US, draw, thick] (not) at (-1.5, -0.2) {};
  \node[or gate US, draw, thick] (or) at (3, 1) {};
  \draw (and1.input 1) -- ++(-3,0) node(ptA){}; \node[left] at (ptA) {$A$};
  \draw (and1.input 2) -- ++(-2.5,0) node(ptB){}; \node[left] at (ptB) {$B$};
  \draw (and1.input 3) -- ++(-2,0) node(ptC){}; \node[left] at (ptC) {$C$};
  \draw (ptA.center) node[circle,fill,inner sep=1pt] {} |- (and2.input 1);
  \draw (ptB.center) node[circle,fill,inner sep=1pt] {} |- (and2.input 2);
  \draw (ptC.center) node[circle,fill,inner sep=1pt] {} |- (not.input);
  \draw (not.output) |- (and2.input 3);
  \draw (and1.output) |- (or.input 1);
  \draw (and2.output) |- (or.input 2);
  \draw (or.output) -- ++(0.5,0) node[right] {$F$};
\end{tikzpicture}
\end{center}


% Ex 26 (Desafio Consenso)
\item \textbf{Desafio:} Simplifique o circuito $F = AB + \overline{A}C + BC$ (Teorema do Consenso) e desenhe sua forma reduzida.
\begin{center}
\begin{tikzpicture}[scale=0.6]
  \node[and gate US, draw, thick] (and1) at (0, 2) {};
  \node[and gate US, draw, thick] (and2) at (0, 0.5) {};
  \node[and gate US, draw, thick] (and3) at (0, -1) {};
  \node[not gate US, draw, thick] (not) at (-1.5, 0.75) {};
  \node[or gate US, draw, thick, logic gate inputs=nnn] (or) at (3.5, 0.5) {};
  \draw (and1.input 1) -- ++(-3.5,0) node(ptA){}; \node[left] at (ptA) {$A$};
  \draw (and1.input 2) -- ++(-1.5,0) node(ptB){}; \node[left] at (ptB) {$B$};
  \draw (and2.input 2) -- ++(-4.5,0) node(ptC){}; \node[left] at (ptC) {$C$};
  \draw (ptA.center) node[circle,fill,inner sep=1pt] {} |- (not.input);
  \draw (not.output) |- (and2.input 1);
  \draw (ptB.center) node[circle,fill,inner sep=1pt] {} |- (and3.input 1);
  \draw (ptC.center) node[circle,fill,inner sep=1pt] {} |- (and3.input 2);
  \draw (and1.output) -- (or.input 1);
  \draw (and2.output) -- (or.input 2);
  \draw (and3.output) -- (or.input 3);
  \draw (or.output) -- ++(0.5,0) node[right] {$F$};
\end{tikzpicture}
\end{center}

\end{enumerate}

\end{document}
